\chapter{Descriptive Statistics and Exploratory Data Analysis}
\label{chap:descriptive-statistics}

Descriptive statistics summarize observed data without introducing a probabilistic model. This distinction should remain explicit: a numerical summary is a property of a dataset, while an estimator is a random variable defined through a sampling model.

\section{Samples and empirical distributions}

Given observations $x_1,\ldots,x_n$, the empirical distribution function is
\begin{align}
F_n(x)=\frac{1}{n}\sum_{i=1}^{n}\mathbf{1}_{\{x_i\leq x\}}.
\end{align}

\section{Location}

The arithmetic mean is
\begin{align}
\bar{x}=\frac{1}{n}\sum_{i=1}^{n}x_i,
\end{align}
while the sample median is an order-statistic-based measure of central location.

\section{Dispersion}

The uncorrected empirical variance is
\begin{align}
m_2=\frac{1}{n}\sum_{i=1}^{n}(x_i-\bar{x})^2,
\end{align}
whereas the conventional unbiased estimator of the population variance under independent sampling is
\begin{align}
s^2=\frac{1}{n-1}\sum_{i=1}^{n}(x_i-\bar{x})^2.
\end{align}
The different denominators reflect different mathematical objects and should not be conflated.

\section{Exploratory graphics}

Histograms, empirical distribution functions, box plots, quantile plots, and scatter plots are treated as exploratory representations. Reproducible plotting code should be stored outside the LaTeX source and exported to `assets/figures/`.
